\documentclass[10pt,a4paper]{amsart}
\usepackage[margin=19mm]{geometry}
\usepackage{amsmath,amssymb,mathtools}
\usepackage[scr=rsfs,cal=euler]{mathalfa}
\usepackage{xfrac}
\usepackage[unicode,hidelinks,bookmarks=false]{hyperref}
\newcommand{\ie}{i.e.\ }
\newcommand{\cf}{cf.\ }
\newcommand{\resp}{respectively\ }
\newcommand{\Subsection}{\subsection}
\providecommand{\nobreakdash}{\nobreak\hbox{-}}
\setlength{\emergencystretch}{2em}
\multlinegap0pt
\usepackage{amssymb}
\usepackage{mathtools}
\usepackage{bm}
\usepackage{dsfont}
\usepackage{float}
\usepackage{tikz}
\usepackage[shortlabels]{enumitem}
\newtheorem{lemma}{Lemma}
\title{\bf\Large Upper Bounds on Tur\'an Densities via Extremal Set Theory}
\author{Yaobin Chen, Xizhi Liu, Ningyuan Yang, Tianming Zhu}
\date{Source: arXiv:2606.02135v1 (CC BY 4.0)}
\begin{document}
\maketitle

\noindent\emph{Source and licence.} \bf\Large Upper Bounds on Tur\'an Densities via Extremal Set Theory. Version arXiv:2606.02135v1 (CC BY 4.0), distributed by arXiv under the Creative Commons Attribution 4.0 International licence, http://creativecommons.org/licenses/by/4.0/, as recorded in the arXiv licence metadata for this exact version and shown on the abstract page https://arxiv.org/abs/2606.02135v1. Full licence notice: \texttt{licenses/arxiv-2606.02135-CC-BY-4.0.txt}.\par\smallskip

\noindent\textit{Editorial scope.} Counted unit: the article's lemma lem:extension-equivalence (source0.tex lines 408--410). The article's complete original proof of it is delivered with it (source0.tex lines 412--414, one \texttt{proof} environment of 3 lines). No mathematics is written, repaired or re-derived: the statement and the proof are the article's own lines, in the article's own notation, and no retained line is substituted or re-flowed.  No further source-local context is needed: every cross-reference of every retained line resolves inside the two delivered spans.  Nothing was omitted from any retained span, so the package carries no omission record.  No retained line contains a citation, so the wrapper carries no bibliography and no bibliography entry.  No retained line contains a figure environment, an image-inclusion command or drawing code, so no image is delivered and the package declares no asset dependency.  Editorial formatting only, in addition: an AMS \texttt{amsart} preamble that restates the article's own declarations and macros used by the retained lines, namely the environments \texttt{lemma} on separate counters, together with the standard packages those lines need.  The retained environments print in this document as Lemma 1 in order of appearance, whereas the article prints the same items with the numbering of its own full document.  The complete native source of the article as distributed in the arXiv e-print for this exact version is bundled unchanged as \texttt{sources/201879/source0.tex}.
\begin{lemma}\label{lem:extension-equivalence}
A partial hypergraph whose partial edges have size at most $r$ admits a homomorphism to an $r$-graph $H$ if and only if its $r$-uniform extension admits a homomorphism to $H$.
\end{lemma}

\begin{proof}
A homomorphism from the extension restricts to a homomorphism from the partial hypergraph.  Conversely, suppose $\phi$ maps every partial edge $e$ injectively into some edge $E_e$ of $H$.  Since $|E_e|=r$, the remaining vertices of $E_e$ can be assigned to the private vertices $U_e$.  These private sets are disjoint and appear in no other partial edge, so the choices are independent and extend $\phi$ to a homomorphism of the extension.
\end{proof}

\end{document}
